\section{Quadro de notação}
\footnotesize
\renewcommand{\arraystretch}{1.1}
\begin{longtable}{@{}p{.28\linewidth}p{.68\linewidth}@{}}
\toprule Símbolo&Significado\\\midrule
\endfirsthead
\toprule Símbolo&Significado\\\midrule
\endhead
\bottomrule\endfoot
$\pi,\pi^{-1}$&Pareamento de Cantor e sua inversa.\\
$Z,Q,Q^*$&Enumerações de $\Z$, $\Q$ e $\Q^*$.\\
$\heightq(p/q)$&Altura $\max\{|p|,q\}$ da fração reduzida.\\
$\Q_H$&Racionais de altura no máximo $H$.\\
$\D,\D_{\rm mid}$&Diádicos e diádicos não inteiros.\\
$J_{k,m},J_{k,m}^{\rm own}$&Célula fechada e célula semiaberta de posse.\\
$x_0,y_0,s$&Centro, valor ancorado e inclinação de um átomo.\\
$\Gamma,\LL,\LL_H$&Gerador aritmético, biblioteca e truncagem por altura.\\
$I_n^x$&Suporte espacial recuperado da interface global.\\
$\Asm(K)$&Montagens contínuas finitas no domínio $K$.\\
$\omega(f;\delta)$&Módulo de continuidade de $f$.\\
$\rho(z)=z_+$&Ativação ReLU.\\
$P_kF,P_Jf,P_\Lambda F$&Interpolantes na malha uniforme, em uma célula e em uma partição.\\
$\EE,\mathbb I_{\Q}$&Gramática de expressões e intervalos de extremos racionais.\\
$P,d$&Precisão transcendental e profundidade de subdivisão intervalar.\\
$M_{\rm int}(J),M(J)$&Majorantes validados de $\|F''\|_{\infty,J}$.\\
$\eta_{\rm cert}(J)$&Certificado local $|J|^2M(J)/8$.\\
$\varepsilon_I,\varepsilon_R$&Tolerâncias de interpolação e racionalização.\\
$H$&Calibre atômico definido em \eqref{eq:46}.\\
\end{longtable}
\normalsize
